% =============================================================================
%  Manuscrito final -- CarniCore -- Entrega 4 (2B) ISR-401, UTEQ
%  Destino: REFSQ 2027, track Research (15 paginas incl. referencias).
%  El track Research Previews (8 paginas) exigiria recortar Trabajo relacionado
%  y Amenazas a la validez; ver resumen_modificaciones.md, apartado 3.
%
%  PLANTILLA OFICIAL: Springer LNCS (llncs.cls).
%  El archivo llncs.cls NO se redistribuye en este repositorio. Descárguelo del
%  sitio oficial de Springer antes de compilar:
%    https://www.springer.com/gp/computer-science/lncs/conference-proceedings-guidelines
%  Coloque llncs.cls junto a este .tex. El estilo splncs04.bst ya está en el
%  repositorio.
%
%  Compilación:
%    pdflatex manuscrito_final && bibtex manuscrito_final && \
%    pdflatex manuscrito_final && pdflatex manuscrito_final
%
%  REGLA DE INTEGRIDAD: ninguna cifra de este documento se escribe a mano.
%  Todas las tablas de resultados se incluyen con \input desde tablas/, que
%  genera 06_Experimento/scripts_analisis/05_generar_figuras.py y
%  06_analisis_potencia.py a partir de los datos crudos. Si una cifra del
%  texto corrido cambia, regenere el pipeline y actualícela desde su salida.
% =============================================================================
\documentclass[runningheads]{llncs}

\usepackage[T1]{fontenc}
\usepackage[utf8]{inputenc}
\usepackage{graphicx}
\usepackage{booktabs}
\usepackage{amsmath}
\usepackage{url}
\usepackage[hidelinks]{hyperref}

\begin{document}

\title{Un detector léxico-sintáctico de ambigüedad no replica el juicio
experto sobre requisitos funcionales en español: estudio de caso en
trazabilidad cárnica}

\titlerunning{Detector léxico frente a juicio experto en RF en español}

\author{%
Ariel O. Castro Baja\~{n}a\orcidID{0009-0005-1575-8935} \and
Kleber O. Crespo Espinoza\orcidID{0009-0000-9145-1357} \and
Edhu X. Gamarra Araujo\orcidID{0009-0001-8312-9656} \and
Carlos A. P\'{e}rez Ruiz\orcidID{0009-0003-6741-9391} \and
Erick J. Quintero Gende\orcidID{0009-0000-6032-4179}}

\authorrunning{A. Castro et al.}

\institute{Universidad T\'{e}cnica Estatal de Quevedo (UTEQ),\\
Facultad de Ciencias de la Computaci\'{o}n, Quevedo, Los R\'{i}os, Ecuador\\
\email{\{acastrob,kcrespoe,egamarraa,cperezr3,equinterog\}@uteq.edu.ec}}

\maketitle

% -----------------------------------------------------------------------------
\begin{abstract}
\textbf{Contexto.} Los detectores automáticos de ambigüedad en requisitos se han
validado casi exclusivamente sobre corpus en inglés. Se desconoce su
comportamiento sobre requisitos funcionales (RF) escritos en español en
dominios agroindustriales latinoamericanos.
\textbf{Objetivo.} Medir en qué grado un detector determinista basado en reglas
léxico-sintácticas replica el juicio de un panel experto sobre los 27~RF de
CarniCore, un sistema real de trazabilidad cárnica en Ecuador.
\textbf{Método.} Estudio de caso con censo completo del corpus y protocolo
pre-registrado en OSF. Un detector de tres categorías (cuantificadores vagos,
conjunciones múltiples y voz pasiva sin agente) se contrastó contra el consenso
por mayoría simple de tres personas expertas que clasificaron los mismos RF de
forma ciega e independiente. Se calcularon $\kappa$ de Cohen por pares,
$\kappa$ de Fleiss, matriz de confusión, precisión, exhaustividad y $F_1$, con
intervalos de confianza al 95\,\% por \textit{bootstrap} de 10.000 réplicas.
\textbf{Resultados.} El detector no marcó ningún requisito (0/27). El panel
marcó cuatro por consenso (RF-08, RF-17, RF-21 y RF-22). La matriz de confusión
resultante es VP\,=\,0, FP\,=\,0, FN\,=\,4, VN\,=\,23, con precisión,
exhaustividad y $F_1$ iguales a 0,0000. El acuerdo interno del panel fue bajo:
$\kappa$ de Fleiss de 0,2636 y $\kappa$ de Cohen por pares entre 0,0870 y
0,3571. La diferencia de prevalencia entre detector y panel no alcanza
significación con este tamaño de corpus (McNemar exacta, $p = 0{,}1250$).
\textbf{Conclusiones.} El resultado nulo tiene dos lecturas simultáneas y ambas
son informativas: el detector léxico no captura la ambigüedad que los expertos
señalan, que es semántica y pragmática; y el propio constructo ``requisito
ambiguo'' resulta poco reproducible entre evaluadores. Publicamos el corpus, las
etiquetas de los tres evaluadores y los guiones de análisis para que el
contraste pueda replicarse.
\keywords{Ambigüedad en requisitos \and Detección automática \and
Requisitos funcionales \and Acuerdo inter-evaluador \and Estudio empírico \and
Ecuador.}
\end{abstract}

% -----------------------------------------------------------------------------
\section{Introducción}
\label{sec:intro}

Las distribuidoras de productos cárnicos en Ecuador gestionan la trazabilidad
animal con registros manuales, lo que dificulta acreditar el cumplimiento
sanitario ante la autoridad competente. CarniCore digitaliza esa cadena
mediante 27~RF redactados en español que cubren registro de proveedores,
ingreso de lotes, cadena de frío, pesaje y consulta pública por
código~QR~\cite{iso29148}, con atributos de calidad especificados según
ISO/IEC~25010:2023~\cite{iso25010}. La calidad léxica de esos requisitos condiciona la
calidad del diseño y de las pruebas: la ambigüedad no detectada propaga
defectos aguas abajo~\cite{mendez2017naming,mendez2015naming}.

La detección automática de ambigüedad se ha estudiado principalmente sobre
corpus en inglés~\cite{ferrari2019nlp,femmer2017smells,gervasi2019ambiguity,%
berry2003ambiguity}. Las herramientas pioneras del área
---ARM~\cite{wilson1997automated} y QuARS~\cite{fabbrini2001linguistic}---
establecieron el enfoque de reglas léxicas que este trabajo evalúa, y
Kiyavitskaya et al.~\cite{kiyavitskaya2008} fijaron los requisitos que debe
cumplir una herramienta de este tipo. La revisión sistemática de Cheng et
al.~\cite{cheng2026genai} sobre inteligencia artificial generativa aplicada a
la ingeniería de requisitos, y el trabajo de Arora et
al.~\cite{arora2024genai}, coinciden en que la evidencia disponible procede de
sistemas de países de renta alta y en inglés. No hemos encontrado estudios que
apliquen un detector a un corpus en español de un dominio agroindustrial
latinoamericano, ni que contrasten esa salida contra un panel experto en ese
contexto.

Un segundo vacío, menos discutido, es el del criterio de referencia. La mayoría
de las evaluaciones de detectores asume que el juicio experto es un patrón oro
estable, pero rara vez reporta el acuerdo interno de ese
panel~\cite{jorgensen2004,shepperd2014researcher}. Si el criterio de referencia
es poco reproducible, las métricas de exactitud diagnóstica que se calculan
contra él heredan esa inestabilidad.

Este trabajo aporta: (1)~un detector de reglas de código abierto para
ambigüedad léxico-sintáctica en español; (2)~la primera medición, hasta donde
alcanza nuestra revisión, del contraste entre ese detector y un panel experto
sobre un corpus real en español del dominio de trazabilidad cárnica,
reportando también el acuerdo interno del panel; y (3)~un paquete de
replicación FAIR~\cite{wilkinson2016fair} con el corpus, las etiquetas
individuales de los tres evaluadores y los guiones de análisis.

\medskip
\noindent\textbf{RQ1.} ¿Con qué exactitud diagnóstica replica el detector
léxico-sintáctico el consenso de un panel experto sobre los 27~RF de CarniCore?

\noindent\textbf{RQ2.} ¿Qué grado de acuerdo alcanzan entre sí las personas
expertas al clasificar esos mismos requisitos, y qué implica ese grado para la
interpretación de RQ1?

% -----------------------------------------------------------------------------
\section{Trabajo relacionado}
\label{sec:related}

Seguimos el esquema abreviado de Kitchenham y Charters~\cite{kitchenham2007},
con las actualizaciones de Petersen et al.~\cite{petersen2015mapping} para la
extracción y clasificación.
La cadena de búsqueda (\texttt{("requirements" AND ("ambiguity" OR "smells")
AND ("detection" OR "NLP"))}) se ejecutó sobre Scopus, IEEE~Xplore y ACM~DL con
ventana 2003--2026; de 148 resultados se revisaron 31 a texto completo y se
retuvieron los 12 más cercanos. La Tabla~\ref{tab:related} resume los seis
principales.

Ferrari y Esuli~\cite{ferrari2019nlp} alcanzan $F_1 \geq 0{,}70$ con modelos de
procesamiento de lenguaje natural para ambigüedad de dominio cruzado en inglés.
Femmer et al.~\cite{femmer2017smells} formalizan doce \textit{smells} de
requisitos detectables por reglas y los aplican a documentación industrial
alemana e inglesa. Gervasi et al.~\cite{gervasi2019ambiguity} y Berry et
al.~\cite{berry2003ambiguity} aportan taxonomías de ambigüedad sin detector
aplicado. Fischbach et al.~\cite{fischbach2023acceptance} combinan
procesamiento de lenguaje natural y modelos de lenguaje para derivar pruebas de
aceptación. Ezzini et al.~\cite{ezzini2022automated} y, antes,
Yang et al.~\cite{yang2011anaphoric} atacan específicamente la ambigüedad
anafórica. Gleich et al.~\cite{gleich2010ambiguity} y Tjong y
Berry~\cite{tjong2013sree} publican detectores de reglas comparables al
nuestro, ambos sobre corpus en inglés y sin contraste contra un panel con
acuerdo reportado.

Ninguno de esos trabajos opera sobre español, ninguno reporta el acuerdo
interno del panel que usa como referencia y ninguno publica el corpus etiquetado
por evaluador individual. Ese es el nicho que ocupa este estudio.

\begin{table}[t]
\centering
\caption{Trabajos más cercanos y diferencia con este estudio.}
\label{tab:related}
\begin{tabular}{p{2.6cm}cc p{2.5cm} p{3.2cm}}
\toprule
\textbf{Referencia} & \textbf{Año} & \textbf{Idioma} & \textbf{Método} &
\textbf{Diferencia con este trabajo} \\
\midrule
Ferrari y Esuli~\cite{ferrari2019nlp} & 2019 & EN & PLN dominio cruzado &
No español; no reporta $\kappa$ del panel \\
Femmer et al.~\cite{femmer2017smells} & 2017 & DE/EN & Reglas (\textit{smells}) &
No español; corpus industrial cerrado \\
Gervasi et al.~\cite{gervasi2019ambiguity} & 2019 & EN & Taxonomía &
Sin detector aplicado \\
Berry et al.~\cite{berry2003ambiguity} & 2003 & EN & Manual de ambigüedad &
Sin evaluación empírica \\
Fischbach et al.~\cite{fischbach2023acceptance} & 2023 & EN & PLN + LLM &
Pruebas de aceptación, no ambigüedad \\
Ezzini et al.~\cite{ezzini2022automated} & 2022 & EN & Anáfora + corpus &
Una sola clase de ambigüedad \\
\textbf{Este trabajo} & \textbf{2026} & \textbf{ES} & \textbf{Reglas + panel} &
\textbf{Español, agroindustria, $\kappa$ del panel publicado} \\
\bottomrule
\end{tabular}
\end{table}

% -----------------------------------------------------------------------------
\section{Método}
\label{sec:method}

\subsection{Diseño y PICOC}

Estudio de caso holístico simple con censo completo del
corpus~\cite{runeson2009,yin2018case}, reportado según la plantilla de
Jedlitschka et al.~\cite{jedlitschka2008} y el cuerpo de conocimiento
SWEBOK~v4.0~\cite{swebok2024}.
\textbf{Población}: los 27~RF verbatim del ERS/SRS~v2.0 de CarniCore
(Pucayacu, La~Man\'{a}, Ecuador).
\textbf{Intervención}: detector determinista basado en reglas.
\textbf{Comparación}: consenso por mayoría simple ($\geq$\,2 de 3) de un panel
de tres personas expertas en ingeniería de requisitos.
\textbf{Resultado}: exactitud diagnóstica del detector frente al consenso, y
acuerdo inter-evaluador del panel.
\textbf{Contexto}: proyecto de fin de curso de pregrado, UTEQ, julio a
septiembre de 2026.

\subsection{Corpus y trabajo de campo}

El corpus procede de un ERS elaborado en tres rondas de campo sobre la empresa
cliente: 16 entrevistas semiestructuradas grabadas en vídeo y audio
(227~minutos acumulados) con siete perfiles operativos, un cuestionario
respondido por 31~personas de seis roles y construido según la lista de
verificación de Molléri et al.~\cite{molleri2020checklist}, cinco documentos
originales de la organización y cinco sesiones de \textit{walkthrough} de
validación. El número de entrevistas se fijó con el criterio de saturación de
significado de Hennink et al.~\cite{hennink2017saturation} y la interpretación
de los resultados se contrastó con participantes del propio estudio conforme a
Birt et al.~\cite{birt2016member}. Todo el
trabajo de campo se ejecutó bajo consentimiento informado conforme a la Ley
Orgánica de Protección de Datos Personales del Ecuador~\cite{lopd2021}; el
material identificable permanece en un contenedor cifrado con AES-256 fuera del
depósito abierto y solo se publican transcripciones seudonimizadas.

\subsection{Detector}

El detector (\texttt{detector\_ambiguedad.py}) aplica tres reglas sobre el texto
verbatim en español de cada RF, derivadas de las categorías de Femmer et
al.~\cite{femmer2017smells} y del manual de Berry et
al.~\cite{berry2003ambiguity}:

\begin{description}
\item[C1 -- Cuantificadores vagos.] Términos cuyo referente no es verificable
sin contexto adicional (\textit{algunos}, \textit{varios}, \textit{adecuado},
\textit{suficiente}, \textit{correspondiente}, entre otros; 26 patrones).
\item[C2 -- Conjunciones múltiples.] Más de tres conectores
\textit{y}/\textit{o}, o más de dos verbos modales en el mismo requisito, como
indicio de requisito compuesto.
\item[C3 -- Voz pasiva sin agente.] Perífrasis pasiva o pasiva refleja sin
agente explícito, que compromete la verificabilidad~\cite{iso29148}.
\end{description}

Un RF se marca como ambiguo si activa al menos una categoría. El análisis es
determinista y reproducible desde el corpus publicado.

\subsection{Panel experto}

Tres personas con experiencia en ingeniería de requisitos clasificaron los
27~RF como ambiguo o no ambiguo de forma ciega e independiente, sin acceso a la
salida del detector, usando una rúbrica común incluida en el material
suplementario. El acuerdo se cuantificó con $\kappa$ de Cohen para cada
par~\cite{cohen1960} y $\kappa$ de Fleiss para el conjunto~\cite{fleiss1971},
con las bandas interpretativas de Landis y Koch~\cite{landis1977}. El consenso
se fijó por mayoría simple antes de conocer los resultados, conforme al
protocolo registrado.

\subsection{Plan de análisis y desviaciones}

El protocolo se registró en OSF (\url{https://osf.io/yp7t3}) el 2 de agosto de
2026, antes de la clasificación experta y de la ejecución del
detector~\cite{nosek2018prereg}. El registro previo es la salvaguarda contra el
\textit{p-hacking}~\cite{simmons2011false} y contra la formulación de hipótesis
después de conocer los resultados~\cite{kerr1998harking}. El plan pre-registrado comprende estadística
descriptiva, matriz de confusión frente al consenso, precisión, exhaustividad y
$F_1$, $\kappa$ de Cohen y de Fleiss, e intervalos de confianza al 95\,\% por
\textit{bootstrap} de 10.000 réplicas con semilla fija~\cite{efron1994},
siguiendo las prácticas de investigación computacional
reproducible~\cite{sandve2013ten,peng2011reproducible}.
Se ejecutó íntegramente. Se registran dos desviaciones:

\begin{description}
\item[DEV-01.] El corpus creció de 25 a 27~RF tras la aprobación de RF-26 y
RF-27 por el comité de control de cambios, con posterioridad al registro. El
detector se aplicó sin modificación al corpus ampliado; los 25~RF originales se
conservan sin alteración y la extensión fue puramente aditiva, lo que se
verificó comparando \texttt{rf25.json} contra \texttt{rf27.json} antes de que
el primero se eliminara al mitigar DEV-03 (ver
\texttt{07\_Datos/desviaciones.md}).
\item[DEV-02.] El protocolo contemplaba una nota interpretativa por cada
desacuerdo entre detector y consenso. Las cuatro notas se redactaron con
posterioridad al análisis y se incluyen como material suplementario, no en el
cuerpo del artículo, por restricción de extensión.
\end{description}

\noindent
Una versión anterior de este manuscrito, circulada internamente el 1 de
septiembre de 2026, declaraba erróneamente que la fase comparativa no se había
ejecutado. Esa declaración era incorrecta: el panel se ejecutó, las métricas se
calcularon y sus salidas estaban versionadas en el repositorio desde esa misma
fecha. Se deja constancia de la corrección por transparencia.

\subsection{Análisis de sensibilidad estadística}

Dado que el estudio es un censo de tamaño fijo, no procede un cálculo de
potencia \textit{a priori} ni una potencia observada, que sería redundante con
el valor $p$~\cite{hoenig2001}. Se reporta en su lugar un análisis de
sensibilidad: qué magnitud de efecto habría sido detectable con
$N = 27$, $\alpha = 0{,}05$ y potencia $0{,}80$, y qué $N$ se requeriría para
efectos de interés (Tabla~\ref{tab:potencia}). El cálculo canónico de
referencia ($d$ de Cohen $=0{,}50$, $\alpha = 0{,}05$, $1-\beta = 0{,}80$)
\cite{cohen1988power,champely2020pwr} se incluye en esa misma tabla.

% -----------------------------------------------------------------------------
\section{Resultados}
\label{sec:results}

\subsection{RQ1: exactitud diagnóstica del detector}

El detector no activó ninguna de las tres categorías sobre ninguno de los
27~RF (Tabla~\ref{tab:detector}, Fig.~\ref{fig:categorias}). El panel marcó
cuatro requisitos por consenso: RF-08 (cálculo de vida útil con cuatro umbrales
por tipo de producto), RF-17 (georreferenciación opcional de la granja de
origen), RF-21 (sincronización de operaciones fuera de línea) y RF-22
(devolución total o parcial con reingreso o baja según estado).

La matriz de confusión resultante (Tabla~\ref{tab:confusion},
Fig.~\ref{fig:porrf}) es VP\,=\,0, FP\,=\,0, FN\,=\,4, VN\,=\,23. Precisión,
exhaustividad y $F_1$ valen 0,0000, con intervalos de confianza al 95\,\% por
\textit{bootstrap} degenerados en $[0{,}0000,\ 0{,}0000]$: ninguna
remuestra de los 27~RF produce un solo acierto, porque el detector no emite
ningún positivo.

Conviene separar dos lecturas. El detector no produce falsos positivos, lo que
en un filtro de primer paso es deseable. Pero tampoco produce verdaderos
positivos, de modo que su valor informativo frente a este panel es nulo: no
aporta ninguna decisión que el equipo no tuviera ya.

\input{tablas/tabla_01_resultados_detector}
\input{tablas/tabla_03_confusion_prf1}

\begin{figure}[t]
\centering
\includegraphics[width=0.78\textwidth]{figuras/figura_01_distribucion_categorias.pdf}
\caption{Activaciones por categoría del detector sobre los 27~RF. Generada por
\texttt{05\_generar\_figuras.py} a partir de
\texttt{dataset\_consolidado.csv}.}
\label{fig:categorias}
\end{figure}

\begin{figure}[t]
\centering
\includegraphics[width=\textwidth]{figuras/figura_02_estado_por_rf.pdf}
\caption{Clasificación por requisito. La serie del detector es plana en cero;
los cuatro positivos corresponden al consenso experto.}
\label{fig:porrf}
\end{figure}

\subsection{RQ2: acuerdo interno del panel}

El acuerdo del panel fue bajo (Tabla~\ref{tab:kappa},
Fig.~\ref{fig:acuerdo}). El $\kappa$ de Fleiss para los tres evaluadores es
0,2636, dentro de la banda \textit{aceptable} de Landis y
Koch~\cite{landis1977} pero en su tercio inferior. Los pares muestran una
dispersión notable: 0,3478 entre los evaluadores 1 y 2, 0,3571 entre 1 y 3, y
0,0870 entre 2 y 3, este último apenas por encima del azar.

\input{tablas/tabla_04_acuerdo_interevaluador}

\begin{figure}[t]
\centering
\includegraphics[width=0.78\textwidth]{figuras/figura_03_acuerdo_expertos.pdf}
\caption{Acuerdo inter-evaluador con las bandas de Landis y Koch.}
\label{fig:acuerdo}
\end{figure}

\subsection{Sensibilidad estadística}

La diferencia de prevalencia entre detector y consenso descansa sobre cuatro
pares discordantes, todos en la misma dirección; la prueba de McNemar exacta
arroja $p = 0{,}1250$, por encima de $\alpha = 0{,}05$. Se necesitan al menos
seis discordantes unidireccionales para alcanzar significación, lo que a la
tasa de discordancia observada equivale a un corpus de unos 41~RF.

El menor $\kappa$ detectable con $N = 27$ y potencia 0,80 es 0,5230; el
$\kappa$ observado (0,2636) queda por debajo de ese umbral, de modo que el
acuerdo del panel no es estadísticamente distinguible del azar con este tamaño
de corpus. Alcanzar un acuerdo moderado ($\kappa = 0{,}40$) con potencia 0,80
exigiría del orden de 47~RF. Por último, la exhaustividad se estima sobre solo
cuatro positivos: aun con cero aciertos, su intervalo exacto al 95\,\% llega
hasta 0,6024.

\input{tablas/tabla_05_potencia}

\begin{figure}[t]
\centering
\includegraphics[width=0.78\textwidth]{figuras/figura_04_curva_potencia.pdf}
\caption{Potencia del contraste $\kappa \neq 0$ frente al tamaño del corpus.}
\label{fig:potencia}
\end{figure}

% -----------------------------------------------------------------------------
\section{Discusión}
\label{sec:discussion}

\subsection{Implicaciones para la investigación}

El resultado admite dos lecturas simultáneas y ninguna de las dos puede
descartarse con estos datos.

La primera es sobre el detector. Los cuatro requisitos que el panel marcó no
contienen ninguno de los marcadores léxicos buscados: RF-08 es ambiguo porque
enumera cuatro umbrales de vida útil sin especificar qué ocurre con tipos de
producto no listados; RF-21 lo es porque no define el comportamiento ante
conflictos de sincronización. Es ambigüedad de alcance y de intención, no de
léxico. Un detector de superficie no puede alcanzarla por diseño, y este
hallazgo delimita empíricamente ese techo sobre un corpus real.

La segunda es sobre el criterio de referencia. Con $\kappa$ de Fleiss de 0,2636
y un par en 0,0870, el propio constructo ``requisito ambiguo'' resulta poco
reproducible entre evaluadores. Esto invita a leer con cautela las métricas de
exactitud publicadas en la literatura del área cuando no van acompañadas del
acuerdo interno de su panel, una omisión que encontramos en los seis trabajos
más cercanos de la Tabla~\ref{tab:related} y que es coherente con lo observado
en otros contextos de juicio experto en ingeniería de
software~\cite{shepperd2014researcher,jorgensen2004}.

\subsection{Implicaciones para la práctica}

Primero, una plantilla de redacción de voz activa con umbrales cuantitativos
explícitos ---la que el equipo aplicó al escribir el ERS--- parece suprimir la
ambigüedad léxica en origen: cero activaciones sobre 27~RF. Si eso se confirma
en réplicas, el esfuerzo de calidad se desplaza de la herramienta al momento de
la escritura.

Segundo, un detector de reglas puede seguir siendo útil como filtro barato de
primer paso sobre corpus heredados o redactados sin plantilla, pero no debe
presentarse como sustituto de la inspección. En este corpus habría dado una
falsa señal de tranquilidad sobre cuatro requisitos que sí requieren
aclaración.

Tercero, para equipos que evalúen detectores: midan y publiquen el acuerdo de
su panel antes de calcular exactitud contra él.

\subsection{Hallazgo contraintuitivo y lecciones aprendidas}

Esperábamos un resultado intermedio y obtuvimos ceros en las tres métricas. La
tentación de presentar el estudio como fallido fue real, y de hecho una versión
preliminar de este manuscrito llegó a omitir la comparación. La lección, para
otros equipos de investigación empírica en pregrado~\cite{fernandez2018napire},
es que un cero bien medido, con su matriz de confusión y su análisis de
sensibilidad, es un resultado; una omisión no lo es. Registrar el protocolo
antes de ver los datos fue lo que hizo evidente la discrepancia.

% -----------------------------------------------------------------------------
\section{Amenazas a la validez}
\label{sec:threats}

Se organizan según Wohlin et al.~\cite{wohlin2012}, con al menos dos amenazas
por categoría, su mitigación y la limitación remanente, atendiendo a la
discusión de Siegmund et al.~\cite{siegmund2015validity} sobre el equilibrio
entre validez interna y externa en ingeniería de software empírica.

\paragraph{Validez interna.}
(i)~\textit{Cobertura del léxico del detector}: los patrones fueron fijados
antes de ver las etiquetas expertas y no se ajustaron después, pero términos
vagos no incluidos en la lista pasarían inadvertidos; mitigado publicando la
lista completa de 26 patrones para su escrutinio.
(ii)~\textit{Contaminación del panel}: los evaluadores podrían haber inferido el
criterio del detector; mitigado entregándoles únicamente el texto de los
requisitos, sin identificadores, prioridad ni fuente, y sin la salida del
detector. Limitación remanente: dos de los tres evaluadores conocían el dominio
del sistema.

\paragraph{Validez de constructo.}
(i)~\textit{Operacionalización de ``ambiguo'' como binario}: la ambigüedad es
gradual y forzarla a dos niveles pierde matiz; mitigado con una rúbrica común,
pero el $\kappa$ bajo sugiere que la definición no fue interpretada de forma
homogénea.
(ii)~\textit{Consenso por mayoría simple}: con tres evaluadores, un solo cambio
de voto altera el consenso de un requisito; se fijó la regla antes del análisis
y se publican las etiquetas individuales para permitir reanálisis con reglas
alternativas.

\paragraph{Validez de conclusión.}
(i)~\textit{Potencia}: con cuatro discordantes, McNemar exacta no alcanza
significación ($p = 0{,}1250$) y el $\kappa$ mínimo detectable (0,5230) supera
al observado; los intervalos se reportan y no se declara significación alguna.
(ii)~\textit{Prevalencia baja}: la exhaustividad se estima sobre cuatro
positivos y su intervalo exacto llega a 0,6024, de modo que no puede
descartarse que un detector moderadamente sensible produjera este resultado por
azar. Ambas limitaciones son intrínsecas al tamaño del corpus.

\paragraph{Validez externa.}
(i)~\textit{Un solo sistema y un solo dominio}: 27~RF de una distribuidora
cárnica ecuatoriana no generalizan a otros dominios ni a otros estilos de
redacción; mitigado publicando el corpus íntegro para réplica.
(ii)~\textit{Efecto de la plantilla de redacción}: el corpus fue escrito con
una plantilla de voz activa que probablemente suprime los patrones buscados, de
modo que el resultado podría no trasladarse a corpus heredados; se declara
explícitamente como condición de contorno del hallazgo.

% -----------------------------------------------------------------------------
\section{Conclusiones y trabajo futuro}
\label{sec:conclusions}

\textbf{RQ1.} El detector léxico-sintáctico no replica el consenso experto sobre
este corpus: 0 de 27 requisitos marcados frente a 4 del panel, con precisión,
exhaustividad y $F_1$ iguales a 0,0000. No produce falsos positivos, pero
tampoco información utilizable.

\textbf{RQ2.} El acuerdo interno del panel es bajo ($\kappa$ de Fleiss 0,2636;
pares entre 0,0870 y 0,3571) y no distinguible del azar con $N = 27$. La
exactitud diagnóstica de RQ1 debe leerse contra un criterio de referencia que
es él mismo inestable.

\textbf{Contribuciones consolidadas.} Se entrega un detector de reglas de código
abierto para español; la primera medición conocida del contraste detector
frente a panel sobre un corpus en español del dominio agroindustrial, con el
acuerdo del panel reportado; y un paquete de replicación FAIR con las etiquetas
individuales de cada evaluador, que ningún trabajo de los revisados publica.

\textbf{Trabajo futuro.} (1)~Ampliar el corpus a los 47~RF que el análisis de
sensibilidad señala como necesarios para detectar un acuerdo moderado.
(2)~Sustituir el detector de reglas por un modelo de lenguaje sobre el mismo
corpus publicado, midiendo si alcanza la ambigüedad semántica que las reglas no
alcanzan. (3)~Replicar el diseño con un panel de al menos cinco evaluadores y
una escala ordinal en lugar de binaria, para comprobar si el acuerdo mejora al
relajar la dicotomía. (4)~Extender el análisis a los requisitos no funcionales
del sistema, en particular a los derivados de la normativa de protección de
datos~\cite{amaral2021gdpr} y a los de explicabilidad de los dos componentes de
aprendizaje automático de
CarniCore~\cite{chazette2020explainability,chazette2022explainable,glinz2007}.

% -----------------------------------------------------------------------------
\section*{Disponibilidad de datos y materiales}

Paquete de replicación (corpus de 27~RF en JSON, etiquetas individuales de los
tres evaluadores, transcripciones seudonimizadas, respuestas del cuestionario,
guiones de análisis en Python y matriz de trazabilidad):
\url{https://doi.org/10.5281/zenodo.22225854}, bajo CC~BY~4.0 para los datos y
MIT para el código, siguiendo los principios FAIR~\cite{wilkinson2016fair} y de
citación de software de Force11~\cite{smith2016software}.
Pre-registro del protocolo: \url{https://osf.io/yp7t3}.
El repositorio declara su cita recomendada en formato Citation File
Format~\cite{druskat2021cff} y su madurez FAIR se autoevaluó con el modelo de
la Research Data Alliance~\cite{rda2020fair}. Código fuente archivado en
Software Heritage~\cite{dicosmo2020referencing,dicosmo2017swh}:
\texttt{swh:1:dir:5741b167af89a201c815b061cc965309d8167069}.
Repositorio: \url{https://github.com/gleiston-guerrero/CarniCore_Proyecto_ISR401}.

\section*{Uso de tecnologías asistidas por inteligencia artificial}

Durante la preparación de este trabajo se utilizó un modelo grande de lenguaje
(Anthropic Claude) para pulir la redacción de párrafos cuyo contenido fue
escrito previamente por las personas autoras, y para revisar el formato
\LaTeX{}, conforme a las políticas editoriales de
Elsevier~\cite{elsevier2023ai} y Springer Nature~\cite{springer2023ai}. Ningún
resultado, cifra, tabla, figura o conclusión fue generado por el modelo: todas
las cifras de este documento proceden de los guiones versionados en
\texttt{06\_Experimento/scripts\_analisis/} ejecutados sobre los datos crudos
publicados. Las personas autoras revisaron el texto resultante y asumen la
responsabilidad íntegra de su contenido.

\section*{Agradecimientos}

Agradecemos a las personas trabajadoras de la distribuidora de cárnicos de
Pucayacu, cantón La~Man\'{a}, provincia de Cotopaxi, que participaron en las
entrevistas y en las sesiones de validación, y a las tres personas evaluadoras
que integraron el panel experto. Trabajo supervisado por el PhD.\ Gleiston
Cicer\'{o}n Guerrero Ulloa, asignatura Ingeniería de Requerimientos (ISR-401),
UTEQ, período 2026--2027~PPA.

\bibliographystyle{splncs04}
\bibliography{referencias}

\end{document}
